% =====================================================================
%  2 Methodology  (largest block — keep ~1 page per half)
%  Source: poster_draft_v2.md ②
% =====================================================================
\section{Methodology}

The methodology is organised into three parts. Section~\ref{sec:hardware} first presents the ESP32-S3-based flight-controller PCB, which provides the hardware foundation on which the firmware operates. Section~\ref{sec:firmware}, which constitutes the core of this work, then details the five-layer firmware architecture. Finally, Section~\ref{sec:setup} describes the experimental setup: static and frequency-domain analyses are performed on the sensor measurements and the estimation algorithms --- the essential components of a flight-control system --- to validate the performance of the overall system. This part also presents the measurement and verification of the real-time determinism of the RTOS task architecture.
% ---------------------------------------------------------------------
%  2.1 Hardware
%  - ESP32-S3 dual-core @ 240 MHz (R14), BLE, native SIMD (PIE via compute HAL)
%  - GY-87: MPU6050 (accel+gyro), HMC5883L (mag), BMP180 (baro)
%  - MicroSD blackbox, WS2812B LEDs, LoRa footprint, power-monitor board
%  - 4-layer KiCad PCB FC_v01 (FC_v02 in progress); gerbers + BOM released (R15)
%  - Figure: PCB render/photo; end with the design-intent sentence
% ---------------------------------------------------------------------
\subsection{Hardware}
\label{sec:hardware}

% ---------------------------------------------------------------------
%  2.2 Firmware  (centre of the paper)
%  - Five-layer general architecture L5 app -> L4 core -> L3 middleware ->
%    L2 drivers -> L1 platform HAL; dependencies point DOWNWARD only (strict
%    downward DAG); no target-specific type leaks into any L2-L4 public header
%    (diagram: FC_v01-general-architecture-v2.html)
%    NOTE: internal engineering audits / decision records are deliberately
%    NOT part of the paper - present the rule itself, not the process.
%  - RTOS task set / periods / core assignment / IPC
%    (diagram: FC_v01-rtos-architecture.html; 14 resident tasks at boot,
%    CALIB_* and BLE_OTA spawn on demand — R13)
%    -> present as Table 1 (task | period | core), data from config.h (R8)
%  - SIMD: hal_simd primitives (dot/add/sub/mul/scale/mat_mul/inv3x3) for
%    float, packed int16, packed int8 -> ESP-DSP (dsps_*/dspm_*) -> S3 PIE
%    SIMD; hal_matrix type-safe facade; zero dynamic allocation;
%    beneficiaries: EKF prediction/update, calibration fits, filtering
% ---------------------------------------------------------------------
\subsection{Firmware}
\label{sec:firmware}

% ---------------------------------------------------------------------
%  2.3 Experimental Setup  (static analyses — how you measured)
%  - Sensor measurements: time-domain histogram/stdev per axis; FFT -> noise
%    floor / interference bands  [NUM: dated log]
%  - Task periods & jitter: nominal from config.h vs SystemView JTAG/TRAX
%    trace (NOT SD); median/p99/max per task, >=1e4 loops, full load + RF
%    [NUM: pending Step 0 re-profile]
%  - Logging robustness: BLE MTU validation (currently never checked), SD
%    fsync latency (3-30 ms) vs ring depth  [NUM: counters + histogram]
%  - Estimation: attitude EKF 7-state quaternion + gyro bias, dual-task async
%    500/100 Hz; altitude EKF 3-state (baro+accel) 100 Hz
%    [NUM: altitude 24.15% is OFFLINE REPLAY ONLY — say so explicitly]
%  - Calibration: gyro static offset; accel 6-DOF ellipsoid (6 orientations);
%    mag 9-DOF ellipsoid; SIMD-accelerated, zero dynamic RAM (keep short)
%  - End with the mandatory protocol note: bench validation, tethered/clamped;
%    free flight NOT yet performed; INDI is future work
% ---------------------------------------------------------------------
\subsection{Experimental Setup}
\label{sec:setup}
